% Figure: Kepler's second law as conserved areal velocity. The radius
% vector sweeps equal areas in equal times, so the body moves fast near
% the focus and slow far from it. Two shaded sweeps of equal area subtend
% a wide arc near periapsis and a narrow arc near apoapsis.
\begin{figure}[htbp]
\centering
\begin{tikzpicture}[lbl/.style={font=\small}]
% ellipse: semi-major a=3.4, semi-minor b=2.4, focus offset c=sqrt(a^2-b^2)
% c = sqrt(11.56-5.76)=sqrt(5.8)=2.408
% the orbit
\draw[engblue, very thick] (0,0) ellipse (3.4 and 2.4);
% the focus (the attracting centre), to the right of geometric centre
\coordinate (F) at (2.408,0);
\fill[engamber] (F) circle (2.6pt);
\node[engamber, lbl, below right] at (F) {focus (centre of force)};
% periapsis sweep (near focus): wide arc on the near side
% near vertex at (a,0)=(3.4,0); sweep a thin wedge of large angular extent
\fill[engred!22]
  (F) -- (3.4,0)
  .. controls (3.25,0.95) and (2.7,1.5) .. (2.0,1.85)
  -- (F);
\draw[engred, thick] (F) -- (3.4,0);
\draw[engred, thick] (F) -- (2.0,1.85);
\node[engred, lbl] at (3.0,0.55) {fast};
% apoapsis sweep (far from focus): narrow arc, equal area, far side
% far vertex at (-a,0)=(-3.4,0)
\fill[englime!28]
  (F) -- (-3.4,0)
  .. controls (-3.3,-0.7) and (-3.05,-1.1) .. (-2.75,-1.45)
  -- (F);
\draw[englime!65!black, thick] (F) -- (-3.4,0);
\draw[englime!65!black, thick] (F) -- (-2.75,-1.45);
\node[englime!55!black, lbl] at (-2.0,-0.55) {slow};
% the body on each arc
\fill[engblack] (3.4,0) circle (2.0pt);
\fill[engblack] (-3.4,0) circle (2.0pt);
\end{tikzpicture}
\caption[Equal areas in equal times]{Kepler's second law as the
conservation of areal velocity. A body under a central force has no
moment about the centre of force, so its angular momentum $L=m
r^2\dot{\theta}$ is constant, and the rate at which the radius vector
sweeps area, $\tfrac{1}{2}r^2\dot{\theta}=L/(2m)$, is constant in time.
The two shaded wedges have equal area and are swept in equal times. Near
the focus the radius is short, so the body must travel through a wide
angle, and a long arc, to sweep that area; far from the focus the radius
is long, so a small angle, and a short arc, suffices. The body therefore
runs fast at periapsis and slow at apoapsis, the qualitative fact that
the areal-velocity law makes quantitative.}
\label{fig:vol03ch04:areal-velocity}
\end{figure}
